\DocumentMetadata{
  pdfversion=2.0,
  pdfstandard=ua-2,
  lang=en-US,
  tagging=on,
  tagging-setup={math/setup=mathml-SE,tabsorder=structure}
}
\documentclass[11pt,letterpaper]{article}
\usepackage[margin=0.68in,includefoot]{geometry}
\usepackage{newcomputermodern}
\usepackage{amsmath,amscd,mathtools}
\usepackage{tikz,adjustbox}
\providecommand{\square}{\mdlgwhtsquare}
\usepackage[hidelinks]{hyperref}
\hypersetup{
  pdftitle={Topology First-Year Exam, May 2021, Part 2},
  pdfauthor={University of Florida Department of Mathematics},
  pdfsubject={Graduate mathematics examination},
  pdfdisplaydoctitle=true,
  bookmarksopen=true
}
\setcounter{secnumdepth}{0}
\setlength{\parindent}{0pt}
\setlength{\parskip}{0.28em}
\setlength{\emergencystretch}{3em}
\setlength{\leftmargini}{2.15em}
\setlength{\leftmarginii}{2.65em}
\setlength{\leftmarginiii}{3em}
\pagestyle{empty}
\raggedbottom
\allowdisplaybreaks
\NewTaggingSocketPlug{block/list/label}{examproblem}{%
  \tagstructbegin{tag=Lbl}\tagstructend%
  \tagstructbegin{tag=\LBody}%
  \tagstructbegin{tag=\ProblemHeadingTag,title-o={\ProblemHeadingText}}%
  \tagmcbegin{tag=\ProblemHeadingTag,actualtext={\ProblemHeadingText}}%
  #2%
  \tagmcend%
  \tagstructend%
}
\newenvironment{examproblems}[2]{%
  \def\ProblemHeadingTag{#2}%
  \AssignTaggingSocketPlug{block/list/label}{examproblem}%
  \begin{enumerate}%
  \setcounter{enumi}{#1}%
  \renewcommand{\labelenumi}{\textbf{\theenumi.}}%
  \setlength{\topsep}{0.35em}%
  \setlength{\partopsep}{0pt}%
  \setlength{\itemsep}{0.48em}%
  \setlength{\parsep}{0.18em}%
}{\end{enumerate}}
\newenvironment{examsubparts}{%
  \AssignTaggingSocketPlug{block/list/label}{default}%
  \begin{enumerate}%
  \setlength{\topsep}{0.18em}%
  \setlength{\partopsep}{0pt}%
  \setlength{\itemsep}{0.16em}%
  \setlength{\parsep}{0.1em}%
}{\end{enumerate}}
\newenvironment{examinstructions}{%
  \par\begingroup\itshape
}{%
  \par\endgroup\vspace{0.18em}
}
\makeatletter
\renewcommand\subsection{\@startsection{subsection}{2}{0pt}%
  {-1.25ex plus -.25ex minus -.1ex}%
  {0.55ex plus .1ex}%
  {\normalfont\large\bfseries}}
\renewcommand{\@maketitle}{%
  \newpage\null\vskip -1em%
  {\footnotesize\bfseries
    \href{https://www.ufl.edu/}{University of Florida}, \href{https://math.ufl.edu/}{Department of Mathematics}\par}%
  \vskip -0.5em%
  \tagstructbegin{tag=H1,title={Topology First-Year Exam, May 2021, Part 2}}%
  \tagpdfparaOff%
  \tagmcbegin{tag=H1}%
    {\LARGE\bfseries\href{https://gma.math.ufl.edu/exams/topology-first-year/}{Topology First-Year Exam}, May 2021, Part 2\par}%
  \tagmcend%
  \tagpdfparaOn%
  \tagstructend%
  \vskip 0.25em%
  \tagmcbegin{artifact}%
    \hrule height 1pt%
  \tagmcend%
  \vskip 1em%
}
\makeatother
\title{Topology First-Year Exam, May 2021, Part 2}
\author{}
\date{}
\begin{document}
\maketitle
\thispagestyle{empty}
\begin{examinstructions}
Answer the following problems and show all your work. Support all statements to the best of your ability.
\end{examinstructions}
\begin{examproblems}{0}{H2}
\def\ProblemHeadingText{Problem 1.}
\item
This problem has two parts.
\begin{examsubparts}
\item[\textnormal{(a)}]
State the Urysohn Lemma.
\item[\textnormal{(b)}]
Let~\(X\) be a connected normal space having two distinct points~\(x\) and~\(y\). Show that~\(X\) has uncountably many points.
\end{examsubparts}
\def\ProblemHeadingText{Problem 2.}
\item
This problem has two parts.
\begin{examsubparts}
\item[\textnormal{(a)}]
State the Brouwer fixed-point theorem for the disc \(B^2\).
\item[\textnormal{(b)}]
Let~\(A\) be a 3 by 3 matrix of positive real numbers. Show that~\(A\) has a positive real eigenvalue. Hint: Let~\(B\) be the intersection of the unit sphere \(S^2\) and the positive octant, \(\{(x,y,z)\in\mathbb{R}^3\mid x,y,z\geq 0\}\).
\end{examsubparts}
\def\ProblemHeadingText{Problem 3.}
\item
Prove that \(\pi_1(S^1,1)\cong\mathbb{Z}\).
\def\ProblemHeadingText{Problem 4.}
\item
This problem has two parts.
\begin{examsubparts}
\item[\textnormal{(a)}]
Define the Hawaiian earring space, \(H\).
\item[\textnormal{(b)}]
Show the images of the generators of the fundamental groups of the circles do not generate the fundamental group of~\(H\).
\end{examsubparts}
\def\ProblemHeadingText{Problem 5.}
\item
This problem has two parts.
\begin{examsubparts}
\item[\textnormal{(a)}]
Define surface.
\item[\textnormal{(b)}]
Compute the fundamental group of the connected sum of a torus and second torus \(T^2\#T^2\).
\end{examsubparts}
\end{examproblems}
\begin{examinstructions}
Answer the following with complete definitions or statements or short proofs.
\end{examinstructions}
\begin{examproblems}{5}{H2}
\def\ProblemHeadingText{Problem 6.}
\item
Show that \(\mathbb{R}^2\) is not homeomorphic to \(\mathbb{R}^3\).
\def\ProblemHeadingText{Problem 7.}
\item
If \(j:A\hookrightarrow X\) is the inclusion of a retract, then show that the induced map on fundamental groups \(j_\#\) is injective.
\def\ProblemHeadingText{Problem 8.}
\item
State the Seifert–van Kampen Theorem.
\def\ProblemHeadingText{Problem 9.}
\item
What is the fundamental group of the torus \(T^2=S^1\times S^1\)?
\def\ProblemHeadingText{Problem 10.}
\item
Given continuous maps of topological spaces \(f:A\to X\) and \(g:A\to Y\), state the universal property satisfied by the pushout.
\end{examproblems}
\end{document}
